Quiero expresar mi profundo agradecimiento a mis directores, el Dr. Renato Paredes Venero y el Dr. Juan Bautista Cabral, por su constante disposición, guía y compromiso a lo largo de todo el desarrollo de este trabajo. Su acompañamiento fue fundamental tanto en los aspectos técnicos como en los momentos de duda o incertidumbre, y su paciencia ante los desafíos inesperados que surgieron durante la reimplementación del modelo resultó invaluable.

Agradezco también a la \unc{} por brindarme una formación académica de excelencia, y en particular a la Facultad de Matemática, Astronomía, Física y Computación por ofrecerme un espacio de estudio, crecimiento y pertenencia durante toda mi carrera. Tanto la institución como su comunidad fueron pilares en mi desarrollo académico y personal.

A mi familia, especialmente a mi madre Fernanda por su amor y apoyo incondicional, a mi padre Alejandro por sus enseñanzas y guía en mi camino y mi hermana Martina por darme fuerzas y recordarme por qué vale la pena siempre seguir adelante. 

A mis hermanos y compañeros de carrera, Guillermo, Lautaro, Tomas, Pedro, Alejo, Victoria y Santiago por las discusiones, viajes, apoyo y por los momentos compartidos que hicieron de estos años una hermosa aventura llena de aprendizajes y risas.

Finalmente, agradezco también a quien se tome el tiempo de leer este trabajo. Espero que encuentre en estas páginas un aporte valioso y bien intencionado, fruto de mucho esfuerzo y dedicación.